The current-STA continuation now tests all five inter-cluster connections on the migrated worst path. A generic selector reads the accepted state's setup report, packed netlist, placement and route; it has no circuit-specific target whitelist. Matching uses the reported driver net and sink cluster location, with ambiguities explicitly retained as a conservative target set rather than an asserted exact atom-pin mapping. The five current targets have unambiguous recorded location matches. The same audited zero-RC, full-device, retained-tree capacity query finds exact path costs equal to their current physical delays: 1.228630, 0.663513, 0.477688, 0.849339 and 0.663513\,ns. No improving candidate exists under that model, so no unchanged candidate is submitted to native analysis. The accepted CPD remains 6.20393\,ns. Frozen initial-target controls also have zero gain on the current state. Both policies exhaust their selected sets (five versus two targets); this is not an equal-work or equal-time performance comparison.

Relaxing other-net occupancy changes only the first target's optimal cost: net 6587 pin 4 decreases from 1.228630 to 0.849339\,ns, a 0.379291\,ns local opportunity. The other four targets have no relaxed fixed-tree improvement. The shorter witness conflicts with three unit-capacity RR nodes owned by nets 2570, 6563 and 7541. A six-query subset experiment releases each singleton and each pair of these networks, preserving all other occupancy and tree conditions; none admits a shorter path. Releasing all three admits the already audited fully relaxed optimum, so no duplicate all-three query is needed. The empty-set result is also reused from the exact fixed-capacity query. Necessity of three releases applies only within this observed three-net set; alternative blocker sets elsewhere in the graph are not excluded.

The conflicting resource of net 2570 serves eight of its nine sinks, while the other two blockers each have one sink. The following intervention therefore examines coordinated repair and shared-subtree handling rather than just refreshing a single target or increasing its query budget. Releasing a network's occupancy is not an actual reroute: all affected sinks must be restored before native FULL/STA acceptance, and blocker timing regressions can erase the target's local gain. This study performs 18 new reference queries (five fixed-capacity, two frozen-target controls, five relaxed, six proper-subset releases) and no new native routing or analysis calls. It rejects simple single-target fixed-tree/free-capacity continuation at this state; it neither establishes global timing optimality nor rules out coordinated, subtree-changing or different-domain interventions.
